%==============================================================================
\section{Khai thác và khám phá (Exploitation and Exploration)}
\label{sec:exploitation-exploration}
%==============================================================================

\begin{ynghia}[title={Hai chiến lược trong RL}]
Trong học máy, \textbf{khám phá (exploration)} là cho phép agent tìm hiểu \textbf{đặc điểm mới} của môi trường;
\textbf{khai thác (exploitation)} là bắt agent \textbf{bám} kiến thức đã có.
Nếu agent chỉ khai thác, dễ \textbf{kẹt} ở giải pháp cục bộ.
Nếu chỉ khám phá, có thể \textbf{không bao giờ} tìm được policy tốt --- đây là \textbf{dilemma khai thác--khám phá}.
\end{ynghia}

\subsection{Khai thác (Exploitation)}

\begin{dinhnghia}[title={Exploitation}]
Khai thác là chiến lược RL trong đó agent dùng \textbf{kiến thức hiện có} để ra quyết định ở một trạng thái nhằm \textbf{tối đa hóa phần thưởng kỳ vọng}.
Mục tiêu: tận dụng những gì đã biết về môi trường để đạt \textbf{kết quả tốt nhất}.
\end{dinhnghia}

\begin{tinhchat}[title={Các khía cạnh chính của khai thác}]
\begin{itemize}
  \item \textbf{Tối đa hóa reward}: chọn hành động theo giá trị và reward đã học để đạt kết quả cao nhất.
  \item \textbf{Hiệu quả ra quyết định}: tập trung hành động reward cao, giảm chi phí tính toán cho khám phá.
  \item \textbf{Quản lý rủi ro}: rủi ro thấp vì ưu tiên hành động đã thử, giảm bất định so với lựa chọn ít quen thuộc.
\end{itemize}
\end{tinhchat}

\subsection{Khám phá (Exploration)}

\begin{dinhnghia}[title={Exploration}]
Khám phá là hành động giúp agent \textbf{thu thập kiến thức} về môi trường hoặc mô hình.
Quá trình khám phá chọn hành động có \textbf{kết quả không chắc chắn} để thu thập thông tin về trạng thái và reward mà hành động sẽ mang lại.
\end{dinhnghia}

\begin{tinhchat}[title={Các khía cạnh chính của khám phá}]
\begin{itemize}
  \item \textbf{Thu thập thông tin}: thử hành động mới ở trạng thái để cải thiện hiểu biết về môi trường hoặc mô hình.
  \item \textbf{Giảm bất định}: giảm độ không chắc chắn về reward và chuyển trạng thái của các hành động chưa thử nhiều.
  \item \textbf{Phủ không gian trạng thái}: với không gian trạng thái lớn hoặc liên tục, khám phá đảm bảo thăm đủ vùng, tránh học \textbf{lệch} vì chỉ dựa trên ít trải nghiệm.
\end{itemize}
\end{tinhchat}

\subsection{Lựa chọn hành động (Action Selection)}

\begin{phuongphap}[title={Mục tiêu huấn luyện agent}]
Mục tiêu RL là dạy agent \textbf{cách hành xử} dưới nhiều trạng thái.
Trong huấn luyện, agent học hành động qua các cách chọn hành động như: \textbf{greedy}, \textbf{epsilon-greedy}, \textbf{Upper Confidence Bound (UCB)}, \ldots
\end{phuongphap}

\subsection{Đánh đổi khai thác và khám phá}

\begin{ynghia}[title={Exploration vs.\ Exploitation trade-off}]
Ý tưởng \textbf{dùng kiến thức hiện có} so với \textbf{thử hành động ngẫu nhiên} gọi là đánh đổi khai thác--khám phá.
Khi khám phá, agent có thể \textbf{cải thiện} kiến thức theo thời gian.
Khi khai thác, agent nhận \textbf{reward cao hơn ngay lập tức}.
Agent không thể làm cả hai tối đa cùng lúc $\Rightarrow$ cần \textbf{cân bằng}.

Phân bổ tài nguyên (thời gian/thử nghiệm) nên phụ thuộc nhu cầu của cả hai hướng, \textbf{luân phiên} theo trạng thái hiện tại và độ phức tạp bài học.
\end{ynghia}

\subsection{Cân bằng khai thác và khám phá}

\begin{phuongphap}[title={Các kỹ thuật cân bằng}]
\begin{itemize}
  \item \textbf{Epsilon-Greedy Action Selection}
  \item \textbf{Multi-Armed Bandit Frameworks}
  \item \textbf{Upper Confidence Bound (UCB)}
\end{itemize}
\end{phuongphap}

\subsubsection{Epsilon-Greedy}

\begin{phuongphap}[title={Epsilon-Greedy Action Selection}]
Trong RL, agent thường chọn hành động theo reward.
Cách tham lam (greedy): luôn chọn hành động \textbf{tối ưu} để tối đa reward ở trạng thái hiện tại.
\textbf{Epsilon-greedy} kết hợp \textbf{khai thác} (dùng kiến thức cũ) và \textbf{khám phá} (tìm lựa chọn mới).
\end{phuongphap}

\begin{phuongphap}[title={Cơ chế epsilon-greedy}]
Phương pháp thường chọn hành động có \textbf{reward kỳ vọng cao nhất}.
Mục tiêu: cân bằng khám phá và khai thác.
Với xác suất nhỏ $\varepsilon$, agent \textbf{khám phá} (chọn ngẫu nhiên hoặc khác greedy) thay vì chỉ khai thác những gì đã học.
\end{phuongphap}

\begin{vidu}[title={Minh họa epsilon-greedy}]
\tsimg{02_exploitation_exploration/img01.jpg}
\end{vidu}

\subsubsection{Multi-Armed Bandit}

\begin{dinhnghia}[title={Multi-Armed Bandit}]
Khung \textbf{multi-armed bandit} cung cấp nền tảng hình thức quản lý cân bằng khai thác--khám phá trong bài toán ra quyết định \textbf{tuần tự}.
Các thuật toán trong khung này phân tích đánh đổi dựa trên \textbf{cơ chế reward} và hoàn cảnh cụ thể (ví dụ mỗi ``cánh tay'' máy đánh bạc cho reward khác nhau).
\end{dinhnghia}

\subsubsection{Upper Confidence Bound (UCB)}

\begin{phuongphap}[title={Upper Confidence Bound}]
UCB là thuật toán phổ biến cân bằng khai thác--khám phá, dựa trên nguyên tắc \textbf{lạc quan trước bất định} (optimism in the face of uncertainty).
UCB chọn hành động tối ưu hóa \textbf{cận trên tin cậy} (upper confidence bound) của reward kỳ vọng: xét cả \textbf{reward trung bình} và \textbf{độ không chắc chắn/biến thiên} của hành động đó.
\end{phuongphap}

\begin{luuy}[title={Gợi ý thực hành}]
$\varepsilon$ lớn $\Rightarrow$ khám phá nhiều, học chậm nhưng tránh kẹt cục bộ.
$\varepsilon$ nhỏ $\Rightarrow$ khai thác mạnh, hội tụ nhanh nhưng dễ bỏ sót hành động tốt hơn.
UCB và bandit phù hợp khi cần cân bằng có \textbf{cơ sở lý thuyết} hơn so với random $\varepsilon$ đơn thuần.
\end{luuy}
